\documentclass[12pt,a4paper]{article}
\usepackage[UTF8,fontset=fandol]{ctex}
\usepackage{amsmath,amssymb,booktabs}
\usepackage[margin=2.5cm]{geometry}
\usepackage{hyperref}

% 与原文一致的宏定义
\newcommand{\method}{DeIR-Dual}
\newcommand{\qbase}{Q_{\text{base}}}
\newcommand{\qpos}{Q_{\text{pos}}}
\newcommand{\qminus}{Q_{\text{minus}}}
\newcommand{\sbase}{S_{\text{base}}}
\newcommand{\spos}{S_{\text{pos}}}
\newcommand{\sneg}{S_{\text{neg}}}
\newcommand{\sfinal}{S_{\text{final}}}

\begin{document}

\title{§4.1--§4.2 中文翻译\\
\large 奖励与惩罚的对称性 \& 零偏置阈值}

\date{}
\maketitle

\section{奖励与惩罚的对称性（$\alpha$ 与 $\beta$）}

为确保指令跟随机制在不同编码器上均能有效且鲁棒地工作，我们超越经验调参，提出了\textbf{尺度对齐原理（Scale Alignment Principle, SAP）}：“指令能量”（无论正向还是负向）应与嵌入空间中的“主题能量”（即基础检索分数）保持量级一致。我们定义 $E_{\text{base}}(d) \equiv S_{\text{base}}(d)$ 为文档在基础通道上的\textbf{信号强度}，并类似地定义 $E_{\text{req}}(d)$ 和 $E_{\text{neg}}(d)$ 分别对应精炼通道与排除通道。SAP 将 $\alpha$ 和 $\beta$ 视为将异构指令信号转换到统一概率尺度的缩放因子——类似于温度缩放（temperature scaling）将各类别的 logit 量级对齐——从而确保基础检索器与我们的双轨修正之间保持\textbf{量级对称性}。

\subsection{惩罚校准（$\alpha$）：中和主题偏置}

我们定义“风险”文档为 $S_{\text{neg}} > \tau$ 的文档（即匹配排除约束的文档）。对于这些文档，惩罚必须足以\textbf{中和其已有的主题动量}——排除文档的力量应等于最初将其检索出来的力量。将 SAP 应用于惩罚通道，得到：

\begin{equation}
  \alpha_{\text{SAP}} =
    \frac{\mathbb{E}[S_{\text{base}} \mid \text{at-risk}]}
         {\mathbb{E}[\text{Softplus}(S_{\text{neg}} - \tau) \mid
         \text{at-risk}]}.
  \label{eq:alpha}
\end{equation}

公式~\eqref{eq:alpha}无需超参数搜索、无需奖励函数、无需任何启发式假设——它仅仅将惩罚校准到与基础信号相匹配的水平。

为验证 $\alpha_{\text{SAP}}$ 并非单一推导策略的偶然产物，我们评估了\textbf{30种独立的推导方法}，横跨六个理论家族（见表~\ref{tab:alpha-derivation}；完整列表见附录~\ref{sec:alpha-appendix}）：

\begin{itemize}
  \item \textbf{A组（主要：尺度对齐）}：期望比值和百分位数匹配方法。$\alpha = 1.000$--$1.030$。
  \item \textbf{B--E组（鲁棒性检验）}：基于分数分辨率统计、分布分离（KL散度）、排序边界分析以及物理启发的半衰期模型。这些方法产生较低的 $\alpha$ 值（$0.05$--$0.50$），过于保守。
  \item \textbf{F组（边界探索）}：基于熵的分辨率、峰度缩放和贝叶斯后验估计。这些方法产生的 $\alpha$ 值高达 $6.15$——数学上有效，但在实践中过于激进，违背了\textit{最小干预}原则。
\end{itemize}

在所有30种策略中，中位数 $\alpha = 0.92$，标准差 $\sigma = 0.15$，而A组方法在标准设置（$\delta = 0.0$）下围绕 $\alpha \approx 1.0$ 展现出\textbf{令人惊讶的共识}。这种跨方法的一致性表明，$\alpha = 1.0$ 不仅仅是一个经验超参数，而是\textbf{一个基本的缩放常数}，它在密集嵌入空间中维护了语义相关性与指令排除之间的量级对称性。

\begin{table}[h]
\centering
\scriptsize
\begin{tabular}{lcccc}
\toprule
\textbf{组别} & \textbf{方法} & $\alpha$ & \textbf{直觉含义} & \textbf{状态} \\
\midrule
A & 尺度对齐 & 1.000 & 期望比值 & \textbf{最优} \\
A & 第50百分位 & 1.000 & 中位数对齐 & 一致 \\
A & 第75百分位 & 1.030 & 75分位对齐 & 接近最优 \\
B & 分数方差 & 0.052 & 编码器分辨率 & 过低 \\
C & KL最小化 & 0.120 & 分布匹配 & 过低 \\
D & NDCG边界 & 0.310 & 排序边界 & 次优 \\
E & 半衰期 & 0.500 & 50\%衰减 & 保守 \\
F & 分数熵 & 6.150 & Shannon分辨率 & 过高 \\
\bottomrule
\end{tabular}
\caption{$\alpha$ 的30种候选推导策略（仅展示代表性子集；完整列表见附录~\ref{sec:alpha-appendix}）。A组方法围绕 $\alpha \approx 1.0$ 展现出惊人的共识。B--E组偏保守（惩罚不足），F组探索了违背最小干预原则的边界情况。}
\label{tab:alpha-derivation}
\end{table}

\subsection{增强校准（$\beta$）：信息溢价}

对于“安全”文档（$S_{\text{neg}} \leq \tau$），增强量级应与基础分数量级相匹配。将 SAP 应用于奖励通道，得到：

\begin{equation}
  \beta = \frac{\mathbb{E}[S_{\text{base}} \mid \text{safe}]}
    {\mathbb{E}[S_{\text{req}} \cdot \text{safety} \mid \text{safe}]}.
  \label{eq:beta}
\end{equation}

从训练集统计量（855条查询，878个正例，12,825个负例）得出 $\beta = 1.926$。$\beta > 1$ 这一事实揭示了一种根本性的\textbf{信息非对称性}：$\alpha$ 仅仅是\textit{中和}已有的主题偏置（逆基础分数而行），而 $\beta$ 必须\textit{注入} $\qpos$ 中蕴含的、$\qbase$ 中不存在的新语义信息。这一\textbf{信息溢价}——即超出 $1.0$ 的额外乘数——量化了将满足特定指令约束的文档“推升”到通用主题背景之上的成本。换言之，$\alpha$ 是阻尼因子，而 $\beta$ 是放大因子，二者的非对称性（$\beta > 1$ vs.\ $\alpha \approx 1$）反映了两个指令通道在本质上的不同。

\section{零偏置阈值（$\delta$）}

阈值偏移量 $\delta$ 控制着排除区域的边界。传统检索模型通常需要人工调谐的噪声裕量（$\delta > 0$）来防止过度抑制。然而，我们的框架采用了一种基于Neyman--Pearson引理的\textbf{零偏置准则}：因为动态语义阈值 $\tau = \cos(\mathbf{h}_{\text{base}}, \mathbf{h}_{\text{neg}})$ 已经内在地包含了每条查询的“语义底噪”，无需额外的偏移量。

更根本地说，\textbf{任何 $\delta \neq 0$ 都会打破 SAP 均衡}。因为 $\alpha$ 和 $\beta$ 是在 $\delta = 0.0$ 的基准下校准的，引入非零偏移会改变风险文档的分布，从而打破公式~\eqref{eq:alpha} 所建立的惩罚量级与基础分数量级之间的对称性。经验上，$\delta > 0$ 会收缩风险文档群体并降低 p-MRR，这证实了\textbf{仅靠动态锚定 $\tau$} 就能提供最优的判别能力，无需额外的人工安全裕量。

\end{document}
